\chapter{Owners and lifetimes}
\label{ch:owners}

\begin{aside}
A signal that outlives its owner is a leak. An effect that
runs after its widget unmounts is a use-after-free. \maya{}'s
ownership model is a tree, like the UI itself.
\end{aside}

\section{The owner tree}

Each signal and effect has an owner --- the lexical scope
that created it. Owners form a tree rooted at the application
context.

When an owner is destroyed (widget unmounted, scope exited),
its owned signals and effects are destroyed too.

\section{Implicit owner}

The current owner is a thread-local set by widget mounting.
Anything created during the widget's mount call is owned by
that widget:

\begin{lstlisting}
ctx.mount(my_widget(), [&] {
    auto count = signal(0);   // owned by my_widget
    ctx.effect([&] {          // owned by my_widget
        log(ctx.get(count));
    });
});
\end{lstlisting}

When \code{my\_widget} unmounts, \code{count} and the
effect are dropped.

\section{Manual owners}

For long-lived state (app-level singletons), use the root
owner:

\begin{lstlisting}
auto user = ctx.root_signal(load_user());
\end{lstlisting}

\code{user} survives until the application context dies.

\section{Detached effects}

Sometimes you want an effect that outlives its caller (a
background sync, for instance). \code{ctx.detached\_effect}
moves ownership to the root.

Use sparingly --- detached effects can leak.

\section{Owner stacks}

Mounting a widget pushes a new owner. Unmounting pops. The
stack is per-thread; the runtime is single-threaded by
design.

\section{Strict mode}

In strict mode, creating a signal outside any owner is an
error. This prevents accidental top-level state that leaks.
\maya{}'s default is strict.

\section{Disposal callbacks}

Some signals carry external resources (file handles, network
connections). Register a disposal callback:

\begin{lstlisting}
auto file = signal_with_dispose(
    open_file("data"),
    [](File* f) { close(f); }
);
\end{lstlisting}

When the owner dies, the callback runs.

\section{Stable references}

A signal handle (\code{SignalHandle<T>}) is a pointer plus
generation counter. Reading from a stale handle (signal
disposed) is a graceful nullopt return, not a crash.

\section{Refcounting alternative}

Some reactive libraries use refcounting (\code{Rc<Signal>}).
\maya{}'s tree model is simpler: no cycles, predictable
destruction order, no atomic refcount overhead.

The trade-off: signals cannot be shared between owner
subtrees without explicit lifting. The lifting helper is
\code{lift\_to\_root}.

\section{Recap}

\begin{recap}
\begin{itemize}[leftmargin=*]
  \item Owners form a tree mirroring the UI.
  \item Mount pushes owner; unmount pops and disposes.
  \item Manual root owners for app-level state.
  \item Detached effects survive their caller; use sparingly.
  \item Disposal callbacks free external resources.
\end{itemize}
\end{recap}

\section*{Workshop}

\begin{enumerate}[leftmargin=*]
  \item Build a widget with internal state. Confirm the
    state is disposed on unmount.
  \item Create a root-owned signal for global config.
  \item Add a disposal callback that closes a file.
\end{enumerate}
